%!TEX root = ../template.tex

%% ----------------------------------------------------------------------------
%%  NOVATHESIS — 3-BackMatter/app-build.tex
%% ----------------------------------------------------------------------------
%%
%% Version 7.11.0 (2026-06-13)
%% Copyright (C) 2004-26 by João M. Lourenço <joao.lourenco@fct.unl.pt>
%%
%% Tutorial appendix: the NOVAthesis build system (Makefile + latexmk +
%% command-line overrides). Intended audience: everyday users building their
%% own document, plus developers and CI users who need to build for specific
%% schools, document types, and languages without editing configuration files.
%%
%% ----------------------------------------------------------------------------

\typeout{NT FILE 3-BackMatter/app-build.tex}%

\chapter{The \ntindex[building!]{Build System}}
\label{app:build}

\novathesis{} documents are built with \ntindex[building!latexmk]{\latexinline{latexmk}}.
All the \hologo{LaTeX}-specific knowledge (default engine, \texttt{biber},
glossaries, auxiliary directory, clean-up lists) lives in the
\verb!latexmkrc! file at the project root, which \texttt{latexmk} loads
automatically.  The \ntindex[building!Makefile]{\latexinline{Makefile}} is a
thin, convenient wrapper around it.

For everyday writing you only need \texttt{make}.  If you cannot (or do not
want to) use \texttt{make}, calling \texttt{latexmk} directly gives the same
result (see Section~\ref{sec:build_nomake}).


%%=======================================================================
\section{Everyday Building}
\label{sec:build_everyday}
%%=======================================================================

\begin{bashcode}
make            # build with LuaLaTeX (recommended default)
make view       # build, then open the PDF
make watch      # rebuild automatically every time a file is saved
make clean      # remove build artifacts (keeps the PDF)
\end{bashcode}

The main targets are:

\begin{center}
\renewcommand{\arraystretch}{1.25}
\begin{ShadedTabular}{lp{8.5cm}}
  \toprule
  \textbf{Target} & \textbf{Effect} \\
  \verb!make! (or \verb!make lua!) & Build with \hologo{LuaLaTeX} \\
  \verb!make pdf!       & Build with \hologo{pdfLaTeX} \\
  \verb!make xe!        & Build with \hologo{XeLaTeX} \\
  \verb!make view!      & Build, then open the \gls{PDF} \\
  \verb!make watch!     & Rebuild on every change (\texttt{latexmk -pvc}) \\
  \verb!make log!       & Show the build log \\
  \verb!make clean!     & Remove build artifacts (keeps the \gls{PDF}) \\
  \verb!make distclean! & \texttt{clean} + remove \gls{PDF} and SyncTeX files \\
  \verb!make help!      & Show the built-in help \\
  \bottomrule
\end{ShadedTabular}
\end{center}

All auxiliary files go into the \verb!AUXDIR/! directory; on success, the
\gls{PDF} (and its SyncTeX file) are copied to the project root.

The behavior may be adjusted with variables appended to the command:

\begin{center}
\renewcommand{\arraystretch}{1.25}
\begin{ShadedTabular}{lp{8.0cm}}
  \toprule
  \textbf{Variable} & \textbf{Effect} \\
  \verb!FILE=name!    & Root \texttt{.tex} file (default: auto-detect the
                        single file containing \latexinline{\documentclass}) \\
  \verb!NT="k=v,..."! & \latexinline{\ntsetup} overrides
                        (see Section~\ref{sec:build_overrides}) \\
  \verb!V=1!          & Verbose: show the raw \hologo{LaTeX} output \\
  \verb!BATCH=1!      & Batch mode, never stops at errors (good for \gls{CI}) \\
  \verb!FASTWRITES=1! & Skip \texttt{morewrites} — much faster
                        \hologo{pdfLaTeX}/\hologo{XeLaTeX} builds
                        (see Section~\ref{sec:build_fastwrites}) \\
  \verb!TL=2024!      & Use \verb!/usr/local/texlive/2024! for this run \\
  \verb!VIEWER=...!   & \gls{PDF} viewer used by \verb!make view! \\
  \verb!FLAGS=...!    & Extra \texttt{latexmk} flags, passed through \\
  \bottomrule
\end{ShadedTabular}
\end{center}

\begin{bashcode}
make V=1                  # see everything LaTeX prints
make pdf BATCH=1          # pdfLaTeX, never stop at errors
make view VIEWER="open -a Skim"
\end{bashcode}


%%=======================================================================
\section{Command-Line Overrides}
\label{sec:build_overrides}
%%=======================================================================

Any \ntindex[configuration!ntsetup]{\latexinline{\ntsetup}} option can be
given at the command line, \emph{without editing}
\ntindex[configuration!0-config/1\_novathesis.tex]{\latexinline{0-Config/1_novathesis.tex}}.
Command-line values take precedence over the values in the configuration
files.

\begin{bashcode}
# Build for a different school/doctype/language, nothing is edited
make NT="doctype=msc,school=uminho/ec,lang=pt"

# Any \ntsetup key works
make NT="docstatus=final,print/index=true"
\end{bashcode}

This is the recommended way to test other schools, to build final versions
in \gls{CI}, or to keep several variants of the same document without
touching the configuration.

Internally, the \verb!NT! variable defines the macro
\latexinline{\ntoverride} before the class loads.  The overrides are applied
right after \verb!0-Config/1_novathesis.tex! is read, and the
\texttt{school} option is then re-applied so that school-derived settings
(e.g.\ the degree acronym, which depends on \texttt{doctype}) stay
consistent, whatever the order of the keys.


%%=======================================================================
\section{Faster \hologo{pdfLaTeX} and \hologo{XeLaTeX} Builds}
\label{sec:build_fastwrites}
%%=======================================================================

\hologo{pdfTeX} and \hologo{XeTeX} can only write to 16 files at a time
(\hologo{LuaTeX} allows 65536).  \novathesis{} alone claims 12 of those 16
slots, so the class loads the \ntindex[packages!morewrites]{\texttt{morewrites}}
package, which removes the limit --- at a steep price: it makes a
\hologo{pdfLaTeX} pass roughly nine times slower.

If your document stays within the limit, you can skip it:

\begin{bashcode}
make pdf FASTWRITES=1
\end{bashcode}

Beyond the 12 slots the class uses, each of the following costs one more,
and only when you actually use it: the symbols glossary (i.e.\ when you
register a symbols file with
\latexinline{\ntaddfile{glossaries}[symbols]{...}} in
\verb!0-Config/4_files.tex!), \latexinline{\listofalgorithms} and
\latexinline{\listoflistings} (when the list is printed), plus whatever your
own packages need (\texttt{minted}, for instance).  The \novathesis{} demo
document uses all of them, so it needs \texttt{morewrites}; most real theses
do not.

If you overshoot, the build stops with
\texttt{No room for a new \textbackslash write} --- simply drop
\verb!FASTWRITES=1! and rebuild.  The switch is a no-op under
\hologo{LuaLaTeX}, which never loads \texttt{morewrites}.

Without \texttt{make}, define the macro before \latexinline{\documentclass}:

\begin{bashcode}
pdflatex -shell-escape -jobname=template \
    '\def\ntnomorewrites{}\input{template}'
\end{bashcode}


%%=======================================================================
\section{Building Without \texttt{make}}
\label{sec:build_nomake}
%%=======================================================================

Since \verb!latexmkrc! carries the configuration, a plain \texttt{latexmk}
call is enough:

\begin{bashcode}
latexmk -pdflua -shell-escape -file-line-error template
\end{bashcode}

The overrides of Section~\ref{sec:build_overrides} are also available:

\begin{bashcode}
# With latexmk
latexmk -pdflua -shell-escape -usepretex \
    -pretex='\def\ntoverride{doctype=msc,lang=pt}' template

# With a plain engine
lualatex -shell-escape -jobname=template \
    '\def\ntoverride{doctype=msc,lang=pt}\input{template}'
\end{bashcode}


%%=======================================================================
\section{Maintainer Tools}
\label{sec:build_maintainer}
%%=======================================================================

The remaining machinery is only relevant for \novathesis{} development and
regression testing.  It lives in \verb!Makefile.dev! (loaded automatically
by the \verb!Makefile! when present; not shipped in releases) and in the
helper script \verb!.Build/nt-variant.sh!.

%%-----------------------------------------------------------------------
\subsection{Building One School Variant}
\label{sec:build_school}
%%-----------------------------------------------------------------------

\begin{bashcode}
make school SCHOOL=nova/fct                     # defaults from schools.conf
make school SCHOOL=uminho/ec TYPE=msc LANG=pt ENGINE=lua
\end{bashcode}

The variant is built with the override mechanism of
Section~\ref{sec:build_overrides}, so the working copy is never modified.
Each variant uses its own auxiliary directory under \verb!AUXDIR/variants/!.
The resulting \gls{PDF} is named
\texttt{univ-school-doctype-lang-engine.pdf} (e.g.\
\verb!nova-fct-phd-en-lua.pdf!).  Defaults come from
\verb!.Build/schools.conf!: the highest degree available for the school, and
\texttt{lua} whenever the school supports it.  If the build fails, the log
file is copied next to where the \gls{PDF} would go, with the same base name.

%%-----------------------------------------------------------------------
\subsection{Building the Full Matrix}
\label{sec:build_matrix}
%%-----------------------------------------------------------------------

\begin{bashcode}
make matrix                  # all variants, sequential
make matrix JOBS=4           # 4 parallel chains
make matrix FILTER=nova        # only NOVA schools
make matrix FILTER=nova/fct    # one school (and its sub-programs)
make matrix FILTER=uminho,ipl  # comma-separated prefixes
\end{bashcode}

\verb!make matrix! builds every school/doctype/language/engine combination
listed in \verb!.Build/schools.conf!.  The \verb!FILTER! variable restricts
the matrix by prefix: a university (\texttt{nova}), a school
(\texttt{nova/fct}, which includes sub-programs such as
\texttt{nova/fct/cbbi}), or a comma-separated list of prefixes.  All \glsfmtshortpl{PDF} of one invocation go
into \verb!AUXDIR/matrix/YYYY-MM-DD@hh-mm-ss! (a folder named with the
invocation timestamp), together with the log/build output of any
failed variants.

The file \verb!schools.conf! is organized in \emph{groups}, each starting
with a header such as:

\begin{verbatim}
[order: school, doctype, lang; processor: pdf,lua]
nova/fct    [phd, msc]  [en, pt]
...
\end{verbatim}

The semantics are:

\begin{itemize}
  \item \textbf{Groups build sequentially}, in the order set by the
        \texttt{order:} tuple inside the group (the last dimension varies
        fastest).  Each variant is built in its \emph{own} auxiliary
        directory: the variants differ only in the command-line
        \latexinline{\ntoverride}, which \texttt{latexmk} does not track as a
        dependency, so a \emph{shared} aux directory would make it skip the
        rebuild and every variant would reuse the first one's \gls{PDF}.
  \item \textbf{Each \texttt{processor:} gets its own chain.}  A group with
        \texttt{processor: pdf,lua} yields two independent sequential chains.
  \item \textbf{\texttt{JOBS} parallelizes chains, never a chain's inside.}
        With \verb!JOBS>1!, group$\times$engine chains run concurrently.
        Concurrent \texttt{biber} processes fail sporadically, so a
        \texttt{flock}-based wrapper (\verb!.Build/shims/biber!) is enabled
        automatically in that case (it can also be forced with
        \verb!BIBER_LOCK=1! for external scripts).
\end{itemize}

%%-----------------------------------------------------------------------
\subsection{Other Maintainer Targets}
\label{sec:build_dev_misc}
%%-----------------------------------------------------------------------

\begin{center}
\renewcommand{\arraystretch}{1.25}
\begin{ShadedTabular}{lp{8.0cm}}
  \toprule
  \textbf{Target} & \textbf{Effect} \\
  \verb!make zip!     & Release archive \texttt{novathesis-<version>@<date>.zip}
                        (excludes \verb!Makefile.dev! and \verb!.Build/!) \\
  \verb!make bump1/2/3! & Bump the major/mid/minor version number \\
  \verb!make needlua! & Regenerate the list of schools that require
                        \hologo{LuaLaTeX} \\
  \verb!make gclean!  & \texttt{git clean -fx} (keeps \verb!Scripts! and
                        \verb!Fonts!) \\
  \verb!make help-dev! & Show the maintainer help \\
  \bottomrule
\end{ShadedTabular}
\end{center}
